% page 68 of the facsimile (image p-067 of the scan)
% block 160.0 x 237.7 mm ; left margin 8.0 mm
\pgc{-12.700mm}{0.000mm}{
\l{\cel{10}60}
\l{}
\l{}
\l{\sou{baterio}. I. (\sou{milito}). Asemblajo de plura kanoni por destruktar}
\l{\cel{3}poziciono enemika. - II. (\sou{Elektr.}) Asemblajo de plura elek-}
\l{\cel{3}tro-kondensatori por produktar granda elektro-descharji.}
\l{\cel{3}- DEFIRS.}
\l{}
\l{\sou{batisto}. Dina stofo, extreme dina. - DEFIRS.}
\l{}
\l{\sou{batrako}. (\sou{zool.})\cel{2}Amfibio, animalo sango-kolda, di qua la tipo}
\l{\cel{3}esas la rano, e qua konstitucas la klaso quaresma di la}
\l{\cel{3}vertebrozi. - DEFIS.}
\l{}
\l{\sou{bavar}.\cel{2}(\sou{netrans.}) Lasar fluar la salivo ek la boko. - FIS.}
\l{}
\l{\sou{bayo}.\cel{2}Mikra golfo di qua la aperturo esas kreta. - DEFIRS.}
\l{}
\l{\sou{bayadero}.\cel{2}Dansistino Indiana. - DEFIRS.}
\l{}
\l{\sou{bayoneto}. Lamo de stalo qua, ajustebla segunvole ye la pinto}
\l{\cel{3}di la fusilo, posibligas uzar olca kom armo di la tipo}
\l{\cel{3}\textquotedbl{}sabro, espado\textquotedbl{}. - DEFIRS.}
\l{}
\l{\sou{bazo}. I. Parto infra di korpo per qua olca su apogas sur to}
\l{\cel{3}quo suportas lu. - II. (\sou{Geom.}) Lineo o plano qua limitizas}
\l{\cel{3}korpo, figuro geometriala, e quan onu prenas kom nivelo}
\l{\cel{3}por mezurar de ibe la alteso-grado di la korpo di ta figuro.}
\l{\cel{3}- DEFIRS.}
\l{}
\l{\sou{bazalto}. (\sou{geol.}) Roko nigratra, de origino fairala, qua ofte}
\l{\cel{3}prizentas fragmenti prisma inter-paralela. - DEFIRS.}
\l{}
\l{\sou{bazaro}. I. Merkato publika, che la Orientani. - II. Edifico}
\l{\cel{3}plur-etajo qua kontenas butiki diversa. III. Loko tektoza,}
\l{\cel{3}galerio ube vendesas objekti omna-sorta. - DEFIRS.}
\l{}
\l{\sou{bazilo}.\cel{2}(\sou{bot.})\cel{2}Genero de planti odoroza, de la familio \textquotedbl{}la}
\l{\cel{3}biacei\textquotedbl{}, di qua un varietato kultivesas en la gardeni kom}
\l{\cel{3}plezur-planto. - L.\cel{1}\sou{ocynum basilicum}. - DEFIRS.}
\l{}
\l{\sou{baziliko}.\cel{2}I. Che la Romani antiqua : edifico rektangula, di-}
\l{\cel{3}vidita aden plura navi da rangi de koloni - uzita kom}
\l{\cel{2}+seanceyo di la tribunali o kom rendevueyo. II. Baziliko}
\l{\cel{3}Romana antiqua, konsakrita a la kulto kristana, e qua ordi-}
\l{\cel{3}nare divenis la kirko precipua di la urbo. - DEFIRS.}
\l{}
\l{\sou{bazilisko}. (\sou{mitol.}) Reptero mitologiala di qua la regardo}
\l{\cel{3}mortigis. - DEFIRS.}
\l{}
\l{\sou{bear}.\cel{2}(\sou{netrans.}) Apertar la boko. (\sou{anke en la senco figurala}:}
\l{\cel{3}\textquotedbl{}truo beanta\textquotedbl{}. - F.}
\l{}
\l{\sou{beata}.\cel{2}Qua ravisesas en Deo. - (\sou{metaf.}) Qua simulas beato ;}
\l{\cel{3}hipokrita. - DEFIRSL.}
}
